% ═══════════════════════════════════════════════════════════════════════════
% MATERIAL-09: Quiz-Karten für Repaso Relámpago
% 10 Fragen für Team-Quiz (Buzzer-System)
% ═══════════════════════════════════════════════════════════════════════════

\documentclass[11pt, a4paper]{article}

\usepackage[utf8]{inputenc}
\usepackage[T1]{fontenc}
\usepackage[ngerman]{babel}
\usepackage{helvet}
\renewcommand{\familydefault}{\sfdefault}

\usepackage[margin=1.5cm]{geometry}
\usepackage[table]{xcolor}
\usepackage{tcolorbox}
\usepackage{array}
\usepackage{tabularx}

% Farben
\definecolor{spanisch}{HTML}{c41e3a}
\definecolor{grau}{HTML}{666666}
\definecolor{hellgrau}{HTML}{f5f5f5}

\tcbuselibrary{skins}

\pagestyle{empty}

\begin{document}

\begin{center}
{\Large\bfseries\textcolor{spanisch}{REPASO RELÁMPAGO -- Quiz-Karten}}\\[0.3cm]
{\normalsize DS 09 | 02.02.2026 | Tag vor der Beobachtung}
\end{center}

\vspace{0.5cm}

\begin{tcolorbox}[
    colback=hellgrau,
    colframe=spanisch,
    title={\bfseries Spielregeln},
    boxrule=0.8pt
]
\begin{itemize}
    \item 2 Teams (gemischt aus Gruppe A + B)
    \item Lehrkraft liest Frage vor
    \item Wer zuerst die Hand hebt, darf antworten
    \item Richtig = 1 Punkt | Falsch = anderes Team darf
    \item \textbf{Antworten auf Spanisch!}
\end{itemize}
\end{tcolorbox}

\vspace{0.5cm}

\renewcommand{\arraystretch}{1.8}

\begin{tabularx}{\textwidth}{|c|X|X|}
\hline
\rowcolor{spanisch!20}
\textbf{\#} & \textbf{Frage} & \textbf{Antwort} \\
\hline
1 & Wie heißt „Bruder" auf Spanisch? & \textit{hermano} \\
\hline
2 & \textit{Soy} oder \textit{estoy} cansado/-a? & \textit{estoy} (Zustand) \\
\hline
3 & Übersetze: „Ich bin 16 Jahre alt." & \textit{Tengo dieciséis años.} \\
\hline
4 & Welches Adjektiv bedeutet „nett, sympathisch"? & \textit{simpático/-a} \\
\hline
5 & Vervollständige: Para ser médico hay que ser... & \textit{paciente / responsable / inteligente...} \\
\hline
6 & Was bedeutet \textit{también}? & auch \\
\hline
7 & \textit{Soy} oder \textit{estoy} inteligente? & \textit{soy} (Eigenschaft) \\
\hline
8 & Wie fragt man: „Wie bist du?" & \textit{¿Cómo eres?} \\
\hline
9 & Nenne 3 Adjektive für eine gute Lehrerin. & z.B. \textit{paciente, inteligente, divertida} \\
\hline
10 & Übersetze: „Mein Vater ist lustig." & \textit{Mi padre es divertido.} \\
\hline
\end{tabularx}

\vspace{1cm}

\begin{center}
{\large\bfseries Bonus-Fragen (bei Gleichstand)}
\end{center}

\vspace{0.3cm}

\begin{tabularx}{\textwidth}{|c|X|X|}
\hline
\rowcolor{spanisch!10}
\textbf{\#} & \textbf{Frage} & \textbf{Antwort} \\
\hline
B1 & Was ist das Gegenteil von \textit{trabajador}? & \textit{perezoso/-a} (faul) \\
\hline
B2 & Nenne 3 Berufe auf Spanisch. & z.B. \textit{médico, profesor, bombero} \\
\hline
B3 & Wie sagt man „Mir gefällt tanzen"? & \textit{Me gusta bailar.} \\
\hline
\end{tabularx}

\vspace{1cm}

\hrule
\vspace{0.3cm}
{\small\color{grau} Hinweis: Einfache Fragen (1, 4, 6) gezielt an schwächere SuS richten. Komplexere Fragen (5, 9) können Team-Beratung erlauben.}

\end{document}
